\section{The Integration Advantage}

\label{sec:integrated}
%% ============================================================================

Individual component comparisons miss the most important advantage of the Hanzo + Lux stack: integration. When 16 components share a common identity model, secrets backend, observability pipeline, and deployment mechanism, the combinatorial cost of integration vanishes.

\subsection{Integration Point Analysis}

A system composed of $n$ independent services has up to $\binom{n}{2} = \frac{n(n-1)}{2}$ pairwise integration points. For $n = 16$ (the number of categories in this paper), this yields 120 potential integration points.

In practice, not all pairs interact. We identify the following integration categories:

\begin{itemize}
    \item \textbf{Auth integration}: Every service must validate identity tokens. With a common IAM, this is configured once. With separate auth providers, each service requires a custom integration. Estimated: 15 integration points.
    \item \textbf{Secrets integration}: Every service needs secrets (API keys, database credentials, TLS certificates). Common KMS eliminates per-service secret management. Estimated: 15 integration points.
    \item \textbf{Observability integration}: Metrics, logs, and traces must be collected from every service. Common format reduces adapter code. Estimated: 15 integration points.
    \item \textbf{Service-to-service communication}: Direct dependencies between services (e.g., inference engine needs vector search, agent framework needs inference engine). Estimated: 25 integration points.
    \item \textbf{Deployment coordination}: Upgrade ordering, health checks, rollback procedures. Estimated: 10 integration points.
\end{itemize}

Total estimated integration points in a point-solution stack: \textbf{80}. In the Hanzo + Lux integrated stack: \textbf{18} (one per K8s CRD, each self-contained).

\subsection{Cost Quantification}

Based on internal engineering data from three production deployments (2024--2026):

\begin{table}[H]
\centering
\caption{Engineering effort comparison (person-months).}
\footnotesize
\begin{tabular}{@{}lrr@{}}
\toprule
Activity & Point solutions & Hanzo + Lux \\
\midrule
Initial setup & 6.0 & 1.5 \\
Auth integration & 3.0 & 0.2 \\
Secrets wiring & 1.5 & 0.1 \\
Observability & 2.0 & 0.3 \\
Service glue code & 4.0 & 0.5 \\
Testing/QA & 3.0 & 1.0 \\
Ongoing maintenance & 2.0/mo & 0.3/mo \\
\midrule
\textbf{Year 1 total} & \textbf{43.5} & \textbf{7.2} \\
\bottomrule
\end{tabular}
\end{table}

The Hanzo + Lux stack requires approximately 6$\times$ less engineering effort in the first year, with the gap widening in subsequent years due to lower maintenance burden.

\subsection{TCO at Scale}

Total cost of ownership for a production AI-blockchain platform serving 1M MAU with 10 inference GPUs:

\begin{table}[H]
\centering
\caption{Annual TCO comparison (USD, thousands).}
\footnotesize
\begin{tabular}{@{}lrr@{}}
\toprule
Cost category & Point solutions & Hanzo + Lux \\
\midrule
SaaS licenses & \$186K & \$0 \\
Compute (K8s) & \$240K & \$240K \\
GPU (inference) & \$180K & \$180K \\
Engineering (ops) & \$360K & \$120K \\
Egress & \$24K & \$0 \\
\midrule
\textbf{Annual total} & \textbf{\$990K} & \textbf{\$540K} \\
\bottomrule
\end{tabular}
\end{table}

SaaS licenses (Auth0 + Pinecone + Mixpanel + Stripe premium + Fireblocks) constitute the largest savings. Engineering operations cost is 3$\times$ lower due to the unified deployment model.

%% ============================================================================
